\section{Rational Functions and Local Rings}

\begin{exercise}
    Let $V = V(Y^2 - X^2(X+1)) \subset \A^2$, and $\overline{X}$, $\overline{Y}$ the residues of $X$, $Y$ in $\Gamma(V)$; let $z = \overline{Y}/\overline{X} \in k(V)$. Find the pole sets of $z$ and of $z^2$. \\
\end{exercise}

\begin{solution}
    First, notice that $z^2 = \overline{X+1}$ so its pole set is the empty set. Let's prove that $(0,0)$ is the only element in the pole set of $z$. Since $z = \overline{Y}/\overline{X}$ and $z = \overline{X(X+1)}/\overline{Y}$, then $(0,0)$ is the only possible element in the pole set. If the pole set is empty, then $z$ must be a rational function, and hence, we have the equality $\overline{Y}/\overline{X} = \overline{p(X,Y)}$ which is equivalent to
    $$Y - Xp(X,Y) = q(X,Y)(Y^2 - X^2(X+1))$$
    in $k[X,Y]$. Taking $X=0$ gives us the equation $Y = q(0,Y)Y^2$ which is impossible. Therefore, the pole set of $z$ is composed of $(0,0)$ only.\\
\end{solution}

\begin{exercise}
    Let $\mathcal{O}_P(V)$ be the local ring of a variety $V$ at a point $P$. Show that there is a natural one-to-one correspondence between the prime ideals in $\mathcal{O}_P(V)$ and the subvarieties of $V$ that pass through $P$. (\textit{Hint:} If $I$ is prime in $\mathcal{O}_P(V)$, $I\cap \Gamma(V)$ is prime in $\Gamma(V)$, and $I$ is generated by $I\cap \Gamma(V)$; use Problem 2.2.)\\
\end{exercise}

\begin{solution}
    By Problem 2.2, we know that there is a natural one-to-one correspondence between the subvarieties of $V$ and the prime ideals of $\Gamma(V)$. If $I$ is a prime ideal of $\mathcal{O}_P(V)$, then $I \cap \Gamma(V)$ is a prime ideal of $\Gamma(V)$. Conversely, if $I$ is a prime ideal of $\Gamma(V)$, then viewed as a subset of $\mathcal{O}_P(V)$, it generates a prime ideal $I'$ of $\mathcal{O}_P(V)$ (same idea as in the proof of Proposition 3). This establishes a natural one-to-one correspondence between the prime ideals of $\Gamma(V)$ and the prime ideals of $\mathcal{O}_P(V)$. Therefore,  there is a natural one-to-one correspondence between the prime ideals in $\mathcal{O}_P(V)$ and the subvarieties of $V$. \\
\end{solution}

\begin{exercise}
    Let $f$ be a rational function on a variety $V$. Let $U = \{P \in V\ | \ f$ is defined at $p\}$. Then $f$ defines a function from $U$ to $k$. Show that this function determines $f$ uniquely. So a rational function may be considered as a type of function, but only on the complement of an algebraic subset of $V$, not on $V$ itself. \\
\end{exercise}

\begin{solution}
    Let $f$ and $g$ be two rational functions in $k(V)$ that have the same induced function, then $h = f-g$ is a rational function in $k(V)$ such that its induced function is the zero function on its domain. Hence, it suffices to show that if $h \in k(V)$ induces the zero function on its domain $U$, then it is the zero rational function. To do so, write $h = \overline{a}/\overline{b}$ where $\overline{a}, \overline{b} \in \Gamma(V)$ and such that $b(P) \neq 0$ for some $P \in V$ (i.e, $h$ is defined on at least one point). By our assumption on $h$, we have that $a(P)/b(P) = 0$ whenever $h$ is defined at $p$, or equivalently, whenever $b(P) \neq 0$. But notice that $a(P)/b(P) = 0$ implies that $a(P) = 0$. Hence, $b(P) \neq 0$ implies that $a(P) = 0$. It follows that the polynomial $ab$ satisfies $a(P)b(P) = 0$ for all $P \in V$. Thus, $ab \in I(V)$, and hence, $\overline{a}\overline{b} = 0$ in $\Gamma(V)$. But since $V$ is a variety, $\Gamma(V)$ is an integral domain so we must have $\overline{a} = 0$ or $\overline{b} = 0$. Since we know that $\overline{b}$ is non-zero on at least one point of $V$, then we get that $\overline{a} = 0$. It follows that $h = 0$ as a rational function. Therefore, $f = g$. \\
\end{solution}

\begin{exercise}
    In the example given in this section, show that it is impossible to write $f = a/b$, where $a,b \in \Gamma(V)$, and $b(P) \neq 0$ for every $P$ where $f$ is defined. Show that the pole set of $f$ is exactly $\{(x,y,z,w) | y=0 \text{ and } w=0\}$. \\
\end{exercise}

\begin{solution}
    First, suppose by contradiction that we can write $f = \overline{a}/\overline{b}$ where
    $$b(x_0,y_0,z_0,w_0) = 0 \implies y_0 = w_0 = 0,$$
    then $b(x_0,y_0,z_0,w_0) \neq 0$ whenever $y_0 \neq 0$ or $w_0 \neq 0$. Hence, if $y_0 \neq 0$, the polynomial $b(x,y_0, z, w)$ has no roots, and hence, it is equal to a non-zero constant. Similarly, when $w_0 \neq 0$, the polynomial $b$ is constant. Let $(x_i,y_i,z_i,w_i) \in \A^4\setminus V(y,w)$ for $i=1,2$. Without loss of generality, we can assume that $y_1, w_2 \neq 0$. Thus, we get that
    $$b(x_1, y_1, z_1, w_1) = b(x_2, y_1, z_2, w_2) = b(x_2, y_2, z_2, w_2) =: c \in \A^1\setminus \{0\}$$
    using our previous remarks. This proves that $b$ is constant on the set $\A^4\setminus V(y,w)$. Now, fix $(x_0,z_0, w_0) \in \A^4$ and consider the polynomial $b_0(y) = b(x_0, y, z_0, w_0)$. By what we proved, we know that this polynomial is constant on the set $\A^1\setminus\{0\}$. Our goal is to show that $b_0(0) = c$, and hence, that $b_0$ is constant over its domain. If $b_0(0) \neq 0$, then $b_0$ has no roots and so it is a constant polynomial. It follows that $b_0 \equiv c$. If $b_0(0) = 0$, then $b_0(y) = y^m b_1(y)$ where $m$ is the multiplicity of the root, and $b_1$ is a polynomial that doesn't vanish at $y = 0$. Since the roots of $b_1$ are necessarily roots of $b_0$, then $b_1$ has no roots. Thus, it is a non-zero constant polynomial, giving us $b_0(y) = Cy^m$ where $C\neq 0$. But notice that $k$ is algebraically closed so there is certainly an $m$th root of 2 which we will call $\alpha$. Thus, $b_0(1) = C$ and $b_0(\alpha) = 2C$, contradicting the fact that $b_0$ is constant on $\A^1\setminus \{0\}$. Hence, the only possible case is that $b_0 \equiv c$. Since it holds for all $(x_0, z_0, w_0) \in \A^3$, we get that $b \equiv c$. Thus, $f = \frac{1}{c}\overline{a} \in \Gamma(V)$. Let's now prove that this is impossible. Notice that $f = \overline{x}/\overline{y}$ implies that
    $$x = \frac{1}{c}ya(x,y,z,w) + h(x,y,z,w)(xw - yz).$$
    Taking $y = w = 0$ gives us $x = 0$, a contradiction. Therefore, by contradiction, such a polynomial $b$ cannot exist.\\

    Next, let's prove that the pole set of $f$ is exactly $V(y,w)$. From the two expressions of $f$: $f = \overline{x}/\overline{y} = \overline{z}/\overline{w}$, we get that the pole set must be a subset of $V(y,w)$. Conversely, let $p = (p_x, 0, p_z, 0) \in V(y,w)$ and suppose by contradiction that we can write $f = \overline{a}/\overline{b}$ where $b(p) \neq 0$. From the equation $\overline{x}/\overline{y} = \overline{a}/\overline{b}$, we get
    $$xb(x,y,z,w) = ya(x,y,z,w) + h_1(x,y,z,w)(xw - yz).$$
    Taking $y = w = 0$ gives us
    $$xb(x,0,z,0) = 0.$$
    But this is impossible because we know that $b(x,0,z,0)$ is a non-zero polynomial ($b(p_x, 0, p_z, 0) \neq 0$). Thus, by contradiction, $p$ must be in the pole set of $f$. It follows that the pole set of $f$ is $V(y,w)$. \\
\end{solution}

\begin{exercise}
    Let $\varphi : V \to W$ be a polynomial map of affine varieties, $\tilde{\varphi} : \Gamma(W) \to \Gamma(V)$ the induced map on coordinate rings. Suppose $P \in V$, $\varphi(P) = Q$. Show that $\tilde{\varphi}$ extends uniquely to a ring homomorphism (also written $\tilde{\varphi}$) from $\mathcal{O}_Q(W)$ to $\mathcal{O}_P(V)$. (Note that $\tilde{\varphi}$ may not extend to all of $k(W)$.) Show that $\tilde{\varphi}(\mathfrak{m}_Q(W)) \subset \mathfrak{m}_P(V)$. \\
\end{exercise}

\begin{solution}
    Define $\psi : \O_Q(W) \to \O_P(V)$ by $\psi(a/b) = \tilde{\varphi}(a)/\tilde{\varphi}(b) = (a \circ \varphi)/(b \circ \varphi)$. First, let's show that its maps to $\O_P(V)$. Let $a/b \in O_Q(W)$ such that $b(Q) \neq 0$, then $(b\circ \varphi)(P) = b(Q) \neq 0$. It follows that $\tilde{\varphi}(a) / \tilde{\varphi}(b) \in O_P(V)$. Next, we need to show that the map is well-defined. Suppose that $a/b = c/d \in O_Q(W)$, then $ad = bc$ as elements of $\Gamma(W)$. Since $\tilde{\varphi}$ is a morphism on $\Gamma(W)$, we get that $\tilde{\varphi}(a)\tilde{\varphi}(d) = \tilde{\varphi}(b)\tilde{\varphi}(c)$, and hence, $\tilde{\varphi}(a)/\tilde{\varphi}(b) = \tilde{\varphi}(c)/\tilde{\varphi}(d)$. Therefore, $\psi$ is a well-defined map from $\O_Q(W)$ to $\O_P(V)$.\\
    
    Now, we need to show that $\psi$ is a ring homomorphism:
    $$\psi(1) = \psi(1/1) = \tilde{\varphi}(1)/\tilde{\varphi}(1) = 1/1 = 1,$$
    and
    \begin{align*}
        \psi((a/b) + (c/d)) &= \psi((ad + bc)/bd) \\
        &= \tilde{\varphi}(ad+bc)/\tilde{\varphi}(bd) \\
        &= (\tilde{\varphi}(a)\tilde{\varphi}(d) + \tilde{\varphi}(b)\tilde{\varphi}(c))/\tilde{\varphi}(b)\tilde{\varphi}(d) \\
        &= (\tilde{\varphi}(a)/\tilde{\varphi}(b)) + (\tilde{\varphi}(c)/\tilde{\varphi}(d)) \\
        &= \psi(a/b) + \psi(c/d),
    \end{align*}
    and 
    \begin{align*}
        \psi((a/b)\cdot (c/d)) &= \psi((ac)/(bd)) \\
        &= \tilde{\varphi}(ac)/\tilde{\varphi}(bd) \\
        &= (\tilde{\varphi}(a)\tilde{\varphi}(c))/(\tilde{\varphi}(b)\tilde{\varphi}(d)) \\
        &= (\tilde{\varphi}(a) / \tilde{\varphi}(b))\cdot (\tilde{\varphi}(c)/\tilde{\varphi}(d)) \\
        &= \psi(a/b)\psi(c/d).\\
    \end{align*}
    Next, we need to show that it extends $\tilde{\varphi}$. For all $a \in \Gamma(W)\subseteq O_Q(W)$, we have
    $$\psi(a) = \psi(a/1) = \tilde{\varphi}(a)/\tilde{\varphi}(1) = \tilde{\varphi}(a)/1 = \tilde{\varphi}(a).$$

    Now, let $\psi_0 : \O_Q(W)\to \O_P(V)$ be a ring homomorphism that extends $\tilde{\varphi}$, then for all $a/b$ in $\O_Q(W)$:
    \begin{align*}
        & \tilde{\varphi}(b)\psi_0(a/b) = \psi_0(b)\psi_0(a/b) = \psi_0(b\cdot (a/b)) = \psi_0(a) = \tilde{\varphi}(a)\\
        &\implies \psi_0(a/b) = \tilde{\varphi}(a)/\tilde{\varphi}(b) = \psi(a/b).
    \end{align*}
    Therefore, $\psi_0 = \psi$, and hence, $\psi$ is unique. Thus, we will now define $\tilde{\varphi}$ to be $\psi$.\\

    Finally, let $a/b \in \m_Q(W)$, then $a(Q) = 0$ which gives us
    $$\tilde{\varphi}(a/b)(P) = (a\circ \varphi)(P)/(b\circ \varphi)(P) = a(Q)/b(Q) = 0$$
    so $\tilde{\varphi}(a/b) \in \m_P(V)$. Thus, $\tilde{\varphi}(\m_Q(W)) \subset \m_P(V)$. \\
\end{solution}

\begin{exercise}
    Let $T : \A^n \to \A^n$ be an affine change of coordinates, $T(P) = Q$. Show that $\tilde{T} : \O_Q(\A^n) \to \O_P(\A^n)$ is an isomorphism. Show that $\tilde{T}$ induces an isomorphism from $\O_Q(V)$ to $\O_P(V^T)$ if $P \in V^T$, for $V$ a subvariety of $\A^n$. \\
\end{exercise}

\begin{solution}
    Since $T$ is an isomoprhism from $\A^n$ to $\A^n$, then $\tilde{T} : \Gamma(\A^n) \to \Gamma(\A^n)$ is also an isomorphism with inverse $\tilde{T}^{-1} : \Gamma(\A^n) \to \Gamma(\A^n)$. By Problem 2.21, both $\tilde{T}$ and $\tilde{T}$ must act by
    $$\tilde{T}(a/b) = \tilde{T}(a)/\tilde{T}(b), \quad \text{ and } \quad \tilde{T}^{-1}(c/d) = \tilde{T}^{-1}(c)/\tilde{T}^{-1}(d).$$
    Hence, for all $a/b \in \O_Q(\A^n)$ and $c/d \in \O_P(\A^n)$:
    $$\tilde{T}^{-1}(\tilde{T}(a/b)) = \tilde{T}^{-1}(\tilde{T}(a))/\tilde{T}^{-1}(\tilde{T}(b)) = a/b$$
    and 
    $$\tilde{T}(\tilde{T}^{-1}(c/d)) = \tilde{T}(\tilde{T}^{-1}(c))/\tilde{T}(\tilde{T}^{-1}(d)) = c/d.$$
    Therefore, $\tilde{T} : \O_Q(\A^n) \to \O_P(\A^n)$ is an isomorphism. The proof that $\tilde{T}$ induces an isomorphism from $\O_Q(V)$ to $\O_P(V^T)$ is the same as the one we just did since $T$ is an isomorphism from $V^T$ to $V$. \\
\end{solution}